% =====================================================================
% Chapter 8: Inside a gap: complex wavenumbers, evanescent waves and the
%            finite plate
% Theory session, dispersion branch (Nov 2026)
% Requires: amsmath, amssymb
% Include in the main file with  % =====================================================================
% Chapter 8: Inside a gap: complex wavenumbers, evanescent waves and the
%            finite plate
% Theory session, dispersion branch (Nov 2026)
% Requires: amsmath, amssymb
% Include in the main file with  % =====================================================================
% Chapter 8: Inside a gap: complex wavenumbers, evanescent waves and the
%            finite plate
% Theory session, dispersion branch (Nov 2026)
% Requires: amsmath, amssymb
% Include in the main file with  % =====================================================================
% Chapter 8: Inside a gap: complex wavenumbers, evanescent waves and the
%            finite plate
% Theory session, dispersion branch (Nov 2026)
% Requires: amsmath, amssymb
% Include in the main file with  \include{chapter_08_complex_k}
% =====================================================================

\chapter{Inside a gap: complex wavenumbers and the finite plate}
\label{ch:complex_k}

The dispersion diagram shows where waves \emph{cannot} propagate. This chapter asks
what the plate does at a frequency inside a gap when a shaker forces it to move
there: the Bloch wavenumber becomes complex, the waves decay from cell to cell, and
the decay rate tells how strongly a finite plate of about ten cells can show the
gap. Numbers come from a Kirchhoff plane-wave model of the production cell
($a = 0.1\,\mathrm{m}$, $h = 5\,\mathrm{mm}$, steel, one resonator per cell,
$f_r = 1200\,\mathrm{Hz}$, $\mu = 0.25$), solved for complex $k$.

\begin{center}
\begin{tabular}{|l|p{0.68\textwidth}|}
\hline
Symbol & Meaning \\
\hline
$k$ & Bloch wavenumber along one direction; here allowed to be complex \\
$k_\mathrm{re}$, $\sigma$ & real and imaginary part of $k$; $\sigma$ is the decay rate in $1/\mathrm{m}$ \\
$\lambda = \omega^2$ & squared angular frequency \\
$\lambda_r$, $\lambda_X$ & resonator eigenvalue $(2\pi\cdot 1200\,\mathrm{Hz})^2$; bare plate eigenvalue at $X$ \\
$\mu$, $\beta = \lambda/\lambda_r$ & mass ratio $0.25$; squared frequency ratio \\
$\mu_\mathrm{eff} = \mu/(1-\beta)$ & effective mass ratio: the resonator acts on the plate like the mass $m/(1-\beta)$ (Chapter~\ref{ch:gap_position}) \\
$D$, $\rho h$ & bending stiffness and mass per area of the plate \\
$\kappa_B$ & bending wavenumber of the bare plate, $\kappa_B^4 = \omega^2\rho h/D$ \\
$p = \pi/a$ & wavenumber at the zone edge $X$ \\
$\delta$ & offset from the zone edge, $k = p + \delta$ \\
$\bar{H}$, $\Delta$, $c$ & mean and half difference of the diagonal, and off-diagonal entry, of the two-level matrix (Chapter~\ref{ch:eigenvectors}) \\
$s = \mathrm{d}\Delta/\mathrm{d}k$ & how fast $\Delta$ changes with $k$ \\
$\eta$ & loss factor of the plate \\
$L$ & distance from source to receiver \\
\hline
\end{tabular}
\end{center}

% ---------------------------------------------------------------------
\section{From $\omega(k)$ to $k(\omega)$}

The dispersion diagram answers: \emph{given a real $k$, which $\omega$ are
possible?} A forced response asks the reverse: the shaker fixes $\omega$, so
\emph{which $k$ are possible?} For real $k$, a frequency inside a gap has no answer.
That does not mean the plate cannot move at that frequency, since the shaker forces
it to. It means that only solutions with \emph{complex} $k$ exist there.

% ---------------------------------------------------------------------
\section{What a complex wavenumber means}

The Bloch theorem states $\psi(x + a) = e^{ika}\psi(x)$. With
$k = k_\mathrm{re} + i\sigma$,
\begin{equation}
  e^{ika} = e^{i k_\mathrm{re} a}\cdot e^{-\sigma a} .
\end{equation}
\begin{itemize}
  \item $e^{i k_\mathrm{re} a}$ is the phase step from one cell to the next, as before;
  \item $e^{-\sigma a}$ is the amplitude factor per cell: each cell multiplies the
        amplitude by $e^{-\sigma a}$.
\end{itemize}
In decibels, the decay per cell is
\begin{equation}
  20\log_{10}\!\left(e^{\sigma a}\right) = 20\,\sigma a\,\log_{10} e = 8.686\,\sigma a
  \ \mathrm{dB}, \qquad \text{or } 8.686\,\sigma \ \mathrm{dB/m} .
\end{equation}

Two special cases:
\begin{itemize}
  \item $k_\mathrm{re} = \pi/a$: $e^{i\pi} = -1$, so neighbouring cells move in
        \emph{antiphase} while the amplitude shrinks: a Bragg-type decaying standing
        pattern;
  \item $k_\mathrm{re} = 0$: neighbouring cells move \emph{in phase} while the amplitude
        shrinks.
\end{itemize}

\paragraph{Roots come in sets of four.} The plate equation contains $k$ only through
squares, so with $k$ also $-k$ is a solution. Without damping its coefficients are
real, so with $k$ also $k^*$ is a solution. The roots therefore come as
$\{k, -k, k^*, -k^*\}$. On each side of a source one keeps the root that decays away
from it.

% ---------------------------------------------------------------------
\section{Evanescent waves exist without resonators}

For the bare Kirchhoff plate, $\kappa^4 = \kappa_B^4$ has the four roots
$\kappa = \pm\kappa_B$ and $\pm i\kappa_B$. The imaginary pair is the \emph{bending
near field}, present at every frequency. At $1100\,\mathrm{Hz}$,
$\kappa_B = 30.1\,\mathrm{m}^{-1}$, so the near field decays by
$\kappa_B a = 3.0$ per cell ($26\,\mathrm{dB}$ per cell). It matters only within
about one cell of a source or an edge.

A gap is therefore not where evanescent waves begin. It is where \emph{the
propagating pair disappears}, and the slowest-decaying remaining wave is what is seen
at a distance.

% ---------------------------------------------------------------------
\section{Two-level model with complex $k$}
\label{sec:two_level_complex}

\subsection{The matrix near $X$}
Near $X$ along $x$, two plate waves matter: $e^{ikx}$ with
$\lambda_1^0 = (D/\rho h)\,k^4$, and its folded partner with
$\lambda_2^0 = (D/\rho h)\,(k - 2\pi/a)^4$. The resonator acts on the plate as a
frequency-dependent mass, so the rigid-mass matrix of Chapter~\ref{ch:two_level}
applies with $\mu$ replaced by $\mu_\mathrm{eff}$:
\begin{equation}
  H_{11} = \lambda_1^0 - \mu_\mathrm{eff}\lambda, \qquad
  H_{22} = \lambda_2^0 - \mu_\mathrm{eff}\lambda, \qquad
  c = -\mu_\mathrm{eff}\lambda .
\end{equation}
As in Chapter~\ref{ch:eigenvectors}, the eigenvalues are
$\lambda = \bar{H} \pm \sqrt{\Delta^2 + c^2}$.

\subsection{Step 1: $\Delta$ and $\bar{H}$ near the zone edge}
Write $k = p + \delta$ with $p = \pi/a$. Then $k - 2\pi/a = \delta - p$, and since
the fourth power is even,
\begin{equation}
  \lambda_1^0 = \frac{D}{\rho h}(p + \delta)^4, \qquad
  \lambda_2^0 = \frac{D}{\rho h}(p - \delta)^4 .
\end{equation}
Expanding both:
\begin{align}
  (p + \delta)^4 &= p^4 + 4p^3\delta + 6p^2\delta^2 + 4p\delta^3 + \delta^4, \\
  (p - \delta)^4 &= p^4 - 4p^3\delta + 6p^2\delta^2 - 4p\delta^3 + \delta^4 .
\end{align}
The difference keeps the odd powers, the sum the even ones:
\begin{align}
  \Delta &= \frac{\lambda_1^0 - \lambda_2^0}{2}
          = \frac{D}{\rho h}\left(4p^3\delta + 4p\delta^3\right)
          = \frac{4\lambda_X}{p}\,\delta\left(1 + \frac{\delta^2}{p^2}\right), \\
  \bar{H} &= \frac{\lambda_1^0 + \lambda_2^0}{2} - \mu_\mathrm{eff}\lambda
          = \lambda_X\left(1 + 6\frac{\delta^2}{p^2} + \frac{\delta^4}{p^4}\right) - \mu_\mathrm{eff}\lambda ,
\end{align}
using $\lambda_X = (D/\rho h)\,p^4$. To first order in $\delta$:
\begin{equation}
  \Delta \approx s\,\delta, \qquad s = \frac{4\lambda_X}{p} = \frac{4\lambda_X a}{\pi},
  \qquad \bar{H} \approx \lambda_X - \mu_\mathrm{eff}\lambda .
\end{equation}
The neglected term $6\delta^2/p^2$ in $\bar{H}$ is about $10\,\%$ in the middle of
this gap.

\subsection{Step 2: insert an imaginary offset}
Set $\delta = i\sigma$. Then $\Delta = is\sigma$ and $\Delta^2 = -s^2\sigma^2$, so
\begin{equation}
  \lambda = \bar{H} \pm \sqrt{c^2 - s^2\sigma^2} .
\end{equation}

\subsection{Step 3: solve for the decay rate}
Squaring, $(\lambda - \bar{H})^2 = c^2 - s^2\sigma^2$, hence
\begin{equation}
  \boxed{\;\sigma = \frac{\sqrt{c^2 - (\lambda - \bar{H})^2}}{s}\;}
  \label{eq:sigma_two_level}
\end{equation}
\begin{itemize}
  \item \emph{Inside the gap}, $|\lambda - \bar{H}| < |c|$: $\sigma$ is real, and
        $k = \pi/a + i\sigma$ is a decaying wave.
  \item \emph{Outside the gap}, the root would be imaginary, which makes $\delta$ real:
        an ordinary propagating wave.
  \item \emph{At an edge}, $\lambda - \bar{H} = \pm|c|$, so $\sigma = 0$. Let
        $\varepsilon$ be the distance of $\lambda$ from that edge. Then
        $c^2 - (\lambda - \bar{H})^2 = (|c| - |\lambda - \bar{H}|)(|c| + |\lambda - \bar{H}|)
        \approx \varepsilon\cdot 2|c|$, so
        \begin{equation}
          \sigma \approx \frac{\sqrt{2|c|\,\varepsilon}}{s} :
        \end{equation}
        $\sigma$ grows like the square root of the distance from the edge.
  \item \emph{With constant coupling}, $\sigma$ is largest at the gap centre:
        $\sigma_\mathrm{max} = |c|/s$, the gap half-width divided by how fast the two
        crossing bands separate in $k$.
\end{itemize}

\paragraph{The same from the band edge.} At a band edge the band is flat:
$\omega \approx \omega_e - \alpha(k - k_e)^2$ with $\alpha > 0$. Solving for $k$,
\begin{equation}
  k - k_e = \pm\sqrt{\frac{\omega_e - \omega}{\alpha}},
\end{equation}
which is real on the band side ($\omega < \omega_e$) and imaginary on the gap side.

\subsection{The resonator version}
Here $c = -\mu_\mathrm{eff}\lambda$ depends on the frequency and grows towards $f_r$.
The decoupled sine wave pins the upper edge at $\lambda_X$:
$\bar{H} = \lambda_X - |c|$. With $x = \lambda_X - \lambda$,
\begin{equation}
  \lambda - \bar{H} = |c| - x
  \quad\Rightarrow\quad
  c^2 - (\lambda - \bar{H})^2 = c^2 - (|c| - x)^2 = x\,(2|c| - x),
\end{equation}
so
\begin{equation}
  \sigma a = \frac{a\sqrt{x(2|c| - x)}}{s} = \frac{\pi}{4}\,\frac{\sqrt{x(2|c| - x)}}{\lambda_X} .
\end{equation}
Worked at $1050\,\mathrm{Hz}$, in units of $\lambda_X$ ($f_X = 1199.8\,\mathrm{Hz}$
in this model):
\begin{align}
  \lambda/\lambda_X &= (1050/1199.8)^2 = 0.7659, \qquad x = 0.2341, \\
  \beta &= (1050/1200)^2 = 0.7656, \qquad \mu_\mathrm{eff} = 0.25/0.2344 = 1.0666, \\
  |c| &= 1.0666\cdot 0.7659 = 0.8169, \\
  x(2|c| - x) &= 0.2341\cdot 1.3997 = 0.3277, \qquad \sigma a = \tfrac{\pi}{4}\sqrt{0.3277} = 0.450 .
\end{align}

\begin{center}
\begin{tabular}{|l|l|l|}
\hline
$f$ in Hz & $\sigma a$, two-wave formula & $\sigma a$, full plane-wave model \\
\hline
900 & 0.236 & 0.276 \\
1050 & 0.450 & 0.407 \\
1150 & 0.528 & 0.346 \\
\hline
\end{tabular}
\end{center}

The two-wave formula gets the size right and degrades towards $f_r$. There
$\mu_\mathrm{eff}$ becomes large, the resonator nearly pins its attachment point, and
many plane waves are needed to describe that, as in Chapter~\ref{ch:pwe_general}.

% ---------------------------------------------------------------------
\section{Results of the full plane-wave model}

\paragraph{Method.} The plane-wave condition of Chapter~\ref{ch:pwe_general},
\begin{equation}
  1 = \mu_\mathrm{eff}\sum_{\mathbf{G}} \frac{\lambda}{\lambda_\mathbf{G}(\mathbf{k}) - \lambda},
  \qquad
  \lambda_\mathbf{G}(\mathbf{k}) = \frac{D}{\rho h}\left[(k_x + G_x)^2 + (k_y + G_y)^2\right]^2,
\end{equation}
with Kirchhoff plate waves and $31\times 31 = 961$ harmonics, is solved at fixed
frequency for complex $k$ along a direction. The square $(k_x + G_x)^2$ is written
without absolute value, so it continues to complex $k$. Roots are located from minima
of the residual on a grid in the complex plane and polished with Newton's method.

\paragraph{Checks.} At real $k$ the model gives the band frequencies $839.2$ and
$1660.6\,\mathrm{Hz}$ at $X$ and $1036.7\,\mathrm{Hz}$ at $M$ (FE: $836$, $1653$,
$1031$). At a pass-band frequency it returns the known real $k$. The complete gap of
this model is $1036.7$--$1199.8\,\mathrm{Hz}$ (FE: $1031$--$1195\,\mathrm{Hz}$).

\paragraph{Along $x$ ($\Gamma$--$X$), slowest-decaying root:}

\begin{center}
\begin{tabular}{|l|l|l|l|}
\hline
$f$ in Hz & $k_\mathrm{re}/(\pi/a)$ & $\sigma a$ & dB/m \\
\hline
800 & 0.921 & 0 & 0 (propagating) \\
841 & 1 & 0.051 & 4.5 \\
850 & 1 & 0.124 & 10.7 \\
900 & 1 & 0.276 & 24.0 \\
1050 & 1 & 0.407 & 35.4 (maximum) \\
1150 & 1 & 0.346 & 30.1 \\
1190 & 1 & 0.191 & 16.6 \\
1198 & 1 & 0.087 & 7.5 \\
\hline
\end{tabular}
\end{center}

\begin{itemize}
  \item The directional gap along $x$ is $839$--$1200\,\mathrm{Hz}$, and the decaying
        root stays at $k_\mathrm{re} = \pi/a$ throughout: antiphase from cell to cell.
  \item Square-root check at the lower edge: $\sigma a/\sqrt{f - 839.2\,\mathrm{Hz}} =
        0.038$ at $841$, $845$ and $850\,\mathrm{Hz}$, constant as predicted.
  \item The near-field root ($k_\mathrm{re} = 0$, $\sigma a \approx 2.9$--$3.7$) is
        always present and plays no role at a distance.
\end{itemize}

\paragraph{Along the diagonal ($\Gamma$--$M$):}

\begin{center}
\begin{tabular}{|l|l|l|l|}
\hline
$f$ in Hz & $k_\mathrm{re}/(\sqrt{2}\pi/a)$ & $\sigma a$ & dB/m \\
\hline
1000 & 0.816 & 0 & 0 \\
1040 & 1 & 0.289 & 25 \\
1050 & 1 & 0.615 & 53 \\
1075 & 0.934 & 1.463 & 127 \\
1150 & 0.697 & 1.532 & 133 \\
1210 & 0.561 & 1.387 & 120 \\
1265 & 0.386 & 0.701 & 61 \\
1300 & 0.549 & 0 & 0 \\
\hline
\end{tabular}
\end{center}

\begin{itemize}
  \item Between $1050$ and $1075\,\mathrm{Hz}$, two roots at the zone edge meet and
        split into a complex pair; after that $k_\mathrm{re}$ leaves the zone edge, and
        the wave oscillates while it decays.
  \item The diagonal gap reaches up to about $1270\,\mathrm{Hz}$ and decays roughly
        four times faster than along $x$.
\end{itemize}

Of the two directions computed, $x$ has the weaker decay everywhere in the complete
gap. Intermediate angles were not computed.

% ---------------------------------------------------------------------
\section{What this means for the finite plate and the insertion loss}

\begin{enumerate}
  \item \emph{Transmission falls roughly like $e^{-\sigma L}$.} For a receiver at
        distance $L$, in-gap transmission drops by about $8.686\,\sigma L$ dB. The
        finite plate spans about ten cells ($1\,\mathrm{m}$); along $x$ at mid-gap that
        is about $35\,\mathrm{dB}$.
  \item \emph{The gap visible on a finite plate is narrower than the Bloch gap.} Near
        the edges $\sigma \to 0$ like a square root. For at least $20\,\mathrm{dB}$ over
        $1\,\mathrm{m}$ one needs $8.686\,\sigma L \ge 20$, i.e.\ $\sigma \ge
        2.3\,\mathrm{m}^{-1}$ or $\sigma a \ge 0.23$. Along $x$ this holds for roughly
        $880$--$1180\,\mathrm{Hz}$, compared with the Bloch gap of
        $1037$--$1200\,\mathrm{Hz}$ in this model.
  \item \emph{Evanescent decay is not damping.} It dissipates no energy: the lattice
        simply cannot carry energy away, so the energy stays near the source or is
        reflected back to it.
  \item \emph{Damping is a separate, weaker decay.} With a loss factor $\eta$ the
        bending wavenumber becomes $\kappa = \kappa_B(1 + i\eta)^{-1/4} \approx
        \kappa_B(1 - i\eta/4)$, so $\sigma = \eta\kappa_B/4$. At $1100\,\mathrm{Hz}$:
        $0.65\,\mathrm{dB/m}$ for $\eta = 0.01$ and $2.6\,\mathrm{dB/m}$ for
        $\eta = 0.04$. Inside the gap the $10$--$35\,\mathrm{dB/m}$ of the evanescent
        decay dominates, except right at the edges, which damping rounds off. Outside
        the gap damping is the only decay and lowers the resonance peaks of the finite
        plate. This is also why undamped Bloch bands suffice for the band positions.
  \item \emph{Few plate modes exist inside the gap.} A mode of the finite plate is a
        superposition of waves bouncing between its edges; inside the gap there is no
        propagating wave to bounce, so only motion localized at the source or at the
        edges remains.
  \item \emph{The driving point is not isolated.} The cell under the shaker still
        moves, through the near field and the decaying wave. The drive-point mobility
        drops much less than the transfer mobility far away. Which insertion-loss
        definition is used (mobility, radiated power) decides how strongly the gap
        shows; to be checked against the forced-response runs.
  \item \emph{Prerequisite:} all of this applies to a finite plate with resonators on
        the $0.1\,\mathrm{m}$ lattice (about 100 small resonators), not to the current
        layouts with 1--6 large resonators.
\end{enumerate}

% ---------------------------------------------------------------------
\section{How FE would compute $k(\omega)$}

In $K_r = T^H K T$ the wavevector enters only through the phase factors
$e^{\pm i k_x a}$ on the periodic images. With $k_y$ fixed and $z = e^{i k_x a}$, the
reduced problem takes the form
\begin{equation}
  \left(\frac{A_{-}}{z} + A_0 + A_{+} z\right)\psi = 0 ,
\end{equation}
with matrices $A$ that contain $\omega^2$: a \emph{quadratic eigenvalue problem} in $z$
at fixed $\omega$. Its roots come in pairs $z$ and $1/z$, i.e.\ $k$ and $-k$. For
complex $k$, $T^H$ has to be replaced by $T(-\mathbf{k})^{\mathsf{T}}$; conjugation is
the same thing only for real $k$. This is the $k(\omega)$ formulation listed as an
optional later step (damped gap depth); it is not needed for the band positions.

% ---------------------------------------------------------------------
\section{Check questions with answers}

\paragraph{Question 1.} \emph{Along $x$ the decaying root has $k_\mathrm{re} = \pi/a$
throughout the gap. What does the motion look like from one cell to the next, and why
must $\sigma$ go to zero at both gap edges?}

\emph{Answer.} Neighbouring cells move in antiphase ($e^{i\pi} = -1$), and the
amplitude shrinks by $e^{-\sigma a}$ per cell; at mid-gap, $\sigma a \approx 0.4$, so
each cell has about $0.67$ of the amplitude of the previous one. Each gap edge is the
flat extreme of a propagating band: the maximum of band 1 at $X$ below, the decoupled
sine wave at $X$ above. There $k(\omega)$ turns from real to complex, and $\sigma$
grows from zero like the square root of the distance from the edge.

\paragraph{Question 2.} \emph{The Bloch complete gap is $1031$--$1195\,\mathrm{Hz}$. On
the $1\,\mathrm{m}$ finite plate, in which part of that range would you expect a clear
insertion loss, and how does a loss factor of $0.04$ change the picture inside and
outside the gap?}

\emph{Answer.} The clearly visible range is narrower than the Bloch gap, because
$\sigma L$ is limited: near the edges $\sigma \to 0$, and $1\,\mathrm{m}$ is not long
enough to build up much attenuation. In the Kirchhoff model, at least $20\,\mathrm{dB}$
over $1\,\mathrm{m}$ along $x$ holds for roughly $880$--$1180\,\mathrm{Hz}$. The gap does
not ``dampen'': evanescent decay dissipates nothing, the lattice just cannot carry the
energy away. A loss factor of $0.04$ adds about $2.6\,\mathrm{dB/m}$ everywhere. Inside
the gap that is small against $10$--$35\,\mathrm{dB/m}$ and only rounds off the edges;
outside the gap it is the only decay and lowers the resonance peaks of the finite
plate, which also enters an insertion loss measured against the bare plate.

\paragraph{Question 3.} \emph{The diagonal decays about four times faster than $x$. A
receiver sits on the finite plate at some angle from the shaker. Which direction
controls how strongly the gap shows up at that receiver, and why?}

\emph{Answer.} Roughly the direction of the line from shaker to receiver relative to
the lattice axes. A receiver along $x$ or $y$ sees the weakest decay (at most about
$35\,\mathrm{dB/m}$), one along the diagonal about four times more. For a quantity
collected over the whole plate, such as radiated power or plate-averaged mobility, the
weakest direction dominates, since most of the energy that gets far travels along $x$
and $y$. Two caveats: the computed rates are for straight wavefronts along $x$ and the
diagonal, not the field of a point source, and on a finite plate edge reflections add
other paths. The faster diagonal decay is not ``four waves being damped'' (there is no
damping in this model, and the factor four is not a count). Two plausible contributions,
not separated: at $M$, four degenerate plate waves form one coupled combination with
$|\varphi|^2 = 4/(\rho h A)$ against $2/(\rho h A)$ at $X$, so the coupling to the
resonator is twice as strong; and $f_r$ lies inside the diagonal gap
($1037$--$1270\,\mathrm{Hz}$), where $\mu_\mathrm{eff}$ is large, while along $x$ it sits
at the upper edge, where the sine wave does not see the resonator at all.

% =====================================================================

\chapter{Inside a gap: complex wavenumbers and the finite plate}
\label{ch:complex_k}

The dispersion diagram shows where waves \emph{cannot} propagate. This chapter asks
what the plate does at a frequency inside a gap when a shaker forces it to move
there: the Bloch wavenumber becomes complex, the waves decay from cell to cell, and
the decay rate tells how strongly a finite plate of about ten cells can show the
gap. Numbers come from a Kirchhoff plane-wave model of the production cell
($a = 0.1\,\mathrm{m}$, $h = 5\,\mathrm{mm}$, steel, one resonator per cell,
$f_r = 1200\,\mathrm{Hz}$, $\mu = 0.25$), solved for complex $k$.

\begin{center}
\begin{tabular}{|l|p{0.68\textwidth}|}
\hline
Symbol & Meaning \\
\hline
$k$ & Bloch wavenumber along one direction; here allowed to be complex \\
$k_\mathrm{re}$, $\sigma$ & real and imaginary part of $k$; $\sigma$ is the decay rate in $1/\mathrm{m}$ \\
$\lambda = \omega^2$ & squared angular frequency \\
$\lambda_r$, $\lambda_X$ & resonator eigenvalue $(2\pi\cdot 1200\,\mathrm{Hz})^2$; bare plate eigenvalue at $X$ \\
$\mu$, $\beta = \lambda/\lambda_r$ & mass ratio $0.25$; squared frequency ratio \\
$\mu_\mathrm{eff} = \mu/(1-\beta)$ & effective mass ratio: the resonator acts on the plate like the mass $m/(1-\beta)$ (Chapter~\ref{ch:gap_position}) \\
$D$, $\rho h$ & bending stiffness and mass per area of the plate \\
$\kappa_B$ & bending wavenumber of the bare plate, $\kappa_B^4 = \omega^2\rho h/D$ \\
$p = \pi/a$ & wavenumber at the zone edge $X$ \\
$\delta$ & offset from the zone edge, $k = p + \delta$ \\
$\bar{H}$, $\Delta$, $c$ & mean and half difference of the diagonal, and off-diagonal entry, of the two-level matrix (Chapter~\ref{ch:eigenvectors}) \\
$s = \mathrm{d}\Delta/\mathrm{d}k$ & how fast $\Delta$ changes with $k$ \\
$\eta$ & loss factor of the plate \\
$L$ & distance from source to receiver \\
\hline
\end{tabular}
\end{center}

% ---------------------------------------------------------------------
\section{From $\omega(k)$ to $k(\omega)$}

The dispersion diagram answers: \emph{given a real $k$, which $\omega$ are
possible?} A forced response asks the reverse: the shaker fixes $\omega$, so
\emph{which $k$ are possible?} For real $k$, a frequency inside a gap has no answer.
That does not mean the plate cannot move at that frequency, since the shaker forces
it to. It means that only solutions with \emph{complex} $k$ exist there.

% ---------------------------------------------------------------------
\section{What a complex wavenumber means}

The Bloch theorem states $\psi(x + a) = e^{ika}\psi(x)$. With
$k = k_\mathrm{re} + i\sigma$,
\begin{equation}
  e^{ika} = e^{i k_\mathrm{re} a}\cdot e^{-\sigma a} .
\end{equation}
\begin{itemize}
  \item $e^{i k_\mathrm{re} a}$ is the phase step from one cell to the next, as before;
  \item $e^{-\sigma a}$ is the amplitude factor per cell: each cell multiplies the
        amplitude by $e^{-\sigma a}$.
\end{itemize}
In decibels, the decay per cell is
\begin{equation}
  20\log_{10}\!\left(e^{\sigma a}\right) = 20\,\sigma a\,\log_{10} e = 8.686\,\sigma a
  \ \mathrm{dB}, \qquad \text{or } 8.686\,\sigma \ \mathrm{dB/m} .
\end{equation}

Two special cases:
\begin{itemize}
  \item $k_\mathrm{re} = \pi/a$: $e^{i\pi} = -1$, so neighbouring cells move in
        \emph{antiphase} while the amplitude shrinks: a Bragg-type decaying standing
        pattern;
  \item $k_\mathrm{re} = 0$: neighbouring cells move \emph{in phase} while the amplitude
        shrinks.
\end{itemize}

\paragraph{Roots come in sets of four.} The plate equation contains $k$ only through
squares, so with $k$ also $-k$ is a solution. Without damping its coefficients are
real, so with $k$ also $k^*$ is a solution. The roots therefore come as
$\{k, -k, k^*, -k^*\}$. On each side of a source one keeps the root that decays away
from it.

% ---------------------------------------------------------------------
\section{Evanescent waves exist without resonators}

For the bare Kirchhoff plate, $\kappa^4 = \kappa_B^4$ has the four roots
$\kappa = \pm\kappa_B$ and $\pm i\kappa_B$. The imaginary pair is the \emph{bending
near field}, present at every frequency. At $1100\,\mathrm{Hz}$,
$\kappa_B = 30.1\,\mathrm{m}^{-1}$, so the near field decays by
$\kappa_B a = 3.0$ per cell ($26\,\mathrm{dB}$ per cell). It matters only within
about one cell of a source or an edge.

A gap is therefore not where evanescent waves begin. It is where \emph{the
propagating pair disappears}, and the slowest-decaying remaining wave is what is seen
at a distance.

% ---------------------------------------------------------------------
\section{Two-level model with complex $k$}
\label{sec:two_level_complex}

\subsection{The matrix near $X$}
Near $X$ along $x$, two plate waves matter: $e^{ikx}$ with
$\lambda_1^0 = (D/\rho h)\,k^4$, and its folded partner with
$\lambda_2^0 = (D/\rho h)\,(k - 2\pi/a)^4$. The resonator acts on the plate as a
frequency-dependent mass, so the rigid-mass matrix of Chapter~\ref{ch:two_level}
applies with $\mu$ replaced by $\mu_\mathrm{eff}$:
\begin{equation}
  H_{11} = \lambda_1^0 - \mu_\mathrm{eff}\lambda, \qquad
  H_{22} = \lambda_2^0 - \mu_\mathrm{eff}\lambda, \qquad
  c = -\mu_\mathrm{eff}\lambda .
\end{equation}
As in Chapter~\ref{ch:eigenvectors}, the eigenvalues are
$\lambda = \bar{H} \pm \sqrt{\Delta^2 + c^2}$.

\subsection{Step 1: $\Delta$ and $\bar{H}$ near the zone edge}
Write $k = p + \delta$ with $p = \pi/a$. Then $k - 2\pi/a = \delta - p$, and since
the fourth power is even,
\begin{equation}
  \lambda_1^0 = \frac{D}{\rho h}(p + \delta)^4, \qquad
  \lambda_2^0 = \frac{D}{\rho h}(p - \delta)^4 .
\end{equation}
Expanding both:
\begin{align}
  (p + \delta)^4 &= p^4 + 4p^3\delta + 6p^2\delta^2 + 4p\delta^3 + \delta^4, \\
  (p - \delta)^4 &= p^4 - 4p^3\delta + 6p^2\delta^2 - 4p\delta^3 + \delta^4 .
\end{align}
The difference keeps the odd powers, the sum the even ones:
\begin{align}
  \Delta &= \frac{\lambda_1^0 - \lambda_2^0}{2}
          = \frac{D}{\rho h}\left(4p^3\delta + 4p\delta^3\right)
          = \frac{4\lambda_X}{p}\,\delta\left(1 + \frac{\delta^2}{p^2}\right), \\
  \bar{H} &= \frac{\lambda_1^0 + \lambda_2^0}{2} - \mu_\mathrm{eff}\lambda
          = \lambda_X\left(1 + 6\frac{\delta^2}{p^2} + \frac{\delta^4}{p^4}\right) - \mu_\mathrm{eff}\lambda ,
\end{align}
using $\lambda_X = (D/\rho h)\,p^4$. To first order in $\delta$:
\begin{equation}
  \Delta \approx s\,\delta, \qquad s = \frac{4\lambda_X}{p} = \frac{4\lambda_X a}{\pi},
  \qquad \bar{H} \approx \lambda_X - \mu_\mathrm{eff}\lambda .
\end{equation}
The neglected term $6\delta^2/p^2$ in $\bar{H}$ is about $10\,\%$ in the middle of
this gap.

\subsection{Step 2: insert an imaginary offset}
Set $\delta = i\sigma$. Then $\Delta = is\sigma$ and $\Delta^2 = -s^2\sigma^2$, so
\begin{equation}
  \lambda = \bar{H} \pm \sqrt{c^2 - s^2\sigma^2} .
\end{equation}

\subsection{Step 3: solve for the decay rate}
Squaring, $(\lambda - \bar{H})^2 = c^2 - s^2\sigma^2$, hence
\begin{equation}
  \boxed{\;\sigma = \frac{\sqrt{c^2 - (\lambda - \bar{H})^2}}{s}\;}
  \label{eq:sigma_two_level}
\end{equation}
\begin{itemize}
  \item \emph{Inside the gap}, $|\lambda - \bar{H}| < |c|$: $\sigma$ is real, and
        $k = \pi/a + i\sigma$ is a decaying wave.
  \item \emph{Outside the gap}, the root would be imaginary, which makes $\delta$ real:
        an ordinary propagating wave.
  \item \emph{At an edge}, $\lambda - \bar{H} = \pm|c|$, so $\sigma = 0$. Let
        $\varepsilon$ be the distance of $\lambda$ from that edge. Then
        $c^2 - (\lambda - \bar{H})^2 = (|c| - |\lambda - \bar{H}|)(|c| + |\lambda - \bar{H}|)
        \approx \varepsilon\cdot 2|c|$, so
        \begin{equation}
          \sigma \approx \frac{\sqrt{2|c|\,\varepsilon}}{s} :
        \end{equation}
        $\sigma$ grows like the square root of the distance from the edge.
  \item \emph{With constant coupling}, $\sigma$ is largest at the gap centre:
        $\sigma_\mathrm{max} = |c|/s$, the gap half-width divided by how fast the two
        crossing bands separate in $k$.
\end{itemize}

\paragraph{The same from the band edge.} At a band edge the band is flat:
$\omega \approx \omega_e - \alpha(k - k_e)^2$ with $\alpha > 0$. Solving for $k$,
\begin{equation}
  k - k_e = \pm\sqrt{\frac{\omega_e - \omega}{\alpha}},
\end{equation}
which is real on the band side ($\omega < \omega_e$) and imaginary on the gap side.

\subsection{The resonator version}
Here $c = -\mu_\mathrm{eff}\lambda$ depends on the frequency and grows towards $f_r$.
The decoupled sine wave pins the upper edge at $\lambda_X$:
$\bar{H} = \lambda_X - |c|$. With $x = \lambda_X - \lambda$,
\begin{equation}
  \lambda - \bar{H} = |c| - x
  \quad\Rightarrow\quad
  c^2 - (\lambda - \bar{H})^2 = c^2 - (|c| - x)^2 = x\,(2|c| - x),
\end{equation}
so
\begin{equation}
  \sigma a = \frac{a\sqrt{x(2|c| - x)}}{s} = \frac{\pi}{4}\,\frac{\sqrt{x(2|c| - x)}}{\lambda_X} .
\end{equation}
Worked at $1050\,\mathrm{Hz}$, in units of $\lambda_X$ ($f_X = 1199.8\,\mathrm{Hz}$
in this model):
\begin{align}
  \lambda/\lambda_X &= (1050/1199.8)^2 = 0.7659, \qquad x = 0.2341, \\
  \beta &= (1050/1200)^2 = 0.7656, \qquad \mu_\mathrm{eff} = 0.25/0.2344 = 1.0666, \\
  |c| &= 1.0666\cdot 0.7659 = 0.8169, \\
  x(2|c| - x) &= 0.2341\cdot 1.3997 = 0.3277, \qquad \sigma a = \tfrac{\pi}{4}\sqrt{0.3277} = 0.450 .
\end{align}

\begin{center}
\begin{tabular}{|l|l|l|}
\hline
$f$ in Hz & $\sigma a$, two-wave formula & $\sigma a$, full plane-wave model \\
\hline
900 & 0.236 & 0.276 \\
1050 & 0.450 & 0.407 \\
1150 & 0.528 & 0.346 \\
\hline
\end{tabular}
\end{center}

The two-wave formula gets the size right and degrades towards $f_r$. There
$\mu_\mathrm{eff}$ becomes large, the resonator nearly pins its attachment point, and
many plane waves are needed to describe that, as in Chapter~\ref{ch:pwe_general}.

% ---------------------------------------------------------------------
\section{Results of the full plane-wave model}

\paragraph{Method.} The plane-wave condition of Chapter~\ref{ch:pwe_general},
\begin{equation}
  1 = \mu_\mathrm{eff}\sum_{\mathbf{G}} \frac{\lambda}{\lambda_\mathbf{G}(\mathbf{k}) - \lambda},
  \qquad
  \lambda_\mathbf{G}(\mathbf{k}) = \frac{D}{\rho h}\left[(k_x + G_x)^2 + (k_y + G_y)^2\right]^2,
\end{equation}
with Kirchhoff plate waves and $31\times 31 = 961$ harmonics, is solved at fixed
frequency for complex $k$ along a direction. The square $(k_x + G_x)^2$ is written
without absolute value, so it continues to complex $k$. Roots are located from minima
of the residual on a grid in the complex plane and polished with Newton's method.

\paragraph{Checks.} At real $k$ the model gives the band frequencies $839.2$ and
$1660.6\,\mathrm{Hz}$ at $X$ and $1036.7\,\mathrm{Hz}$ at $M$ (FE: $836$, $1653$,
$1031$). At a pass-band frequency it returns the known real $k$. The complete gap of
this model is $1036.7$--$1199.8\,\mathrm{Hz}$ (FE: $1031$--$1195\,\mathrm{Hz}$).

\paragraph{Along $x$ ($\Gamma$--$X$), slowest-decaying root:}

\begin{center}
\begin{tabular}{|l|l|l|l|}
\hline
$f$ in Hz & $k_\mathrm{re}/(\pi/a)$ & $\sigma a$ & dB/m \\
\hline
800 & 0.921 & 0 & 0 (propagating) \\
841 & 1 & 0.051 & 4.5 \\
850 & 1 & 0.124 & 10.7 \\
900 & 1 & 0.276 & 24.0 \\
1050 & 1 & 0.407 & 35.4 (maximum) \\
1150 & 1 & 0.346 & 30.1 \\
1190 & 1 & 0.191 & 16.6 \\
1198 & 1 & 0.087 & 7.5 \\
\hline
\end{tabular}
\end{center}

\begin{itemize}
  \item The directional gap along $x$ is $839$--$1200\,\mathrm{Hz}$, and the decaying
        root stays at $k_\mathrm{re} = \pi/a$ throughout: antiphase from cell to cell.
  \item Square-root check at the lower edge: $\sigma a/\sqrt{f - 839.2\,\mathrm{Hz}} =
        0.038$ at $841$, $845$ and $850\,\mathrm{Hz}$, constant as predicted.
  \item The near-field root ($k_\mathrm{re} = 0$, $\sigma a \approx 2.9$--$3.7$) is
        always present and plays no role at a distance.
\end{itemize}

\paragraph{Along the diagonal ($\Gamma$--$M$):}

\begin{center}
\begin{tabular}{|l|l|l|l|}
\hline
$f$ in Hz & $k_\mathrm{re}/(\sqrt{2}\pi/a)$ & $\sigma a$ & dB/m \\
\hline
1000 & 0.816 & 0 & 0 \\
1040 & 1 & 0.289 & 25 \\
1050 & 1 & 0.615 & 53 \\
1075 & 0.934 & 1.463 & 127 \\
1150 & 0.697 & 1.532 & 133 \\
1210 & 0.561 & 1.387 & 120 \\
1265 & 0.386 & 0.701 & 61 \\
1300 & 0.549 & 0 & 0 \\
\hline
\end{tabular}
\end{center}

\begin{itemize}
  \item Between $1050$ and $1075\,\mathrm{Hz}$, two roots at the zone edge meet and
        split into a complex pair; after that $k_\mathrm{re}$ leaves the zone edge, and
        the wave oscillates while it decays.
  \item The diagonal gap reaches up to about $1270\,\mathrm{Hz}$ and decays roughly
        four times faster than along $x$.
\end{itemize}

Of the two directions computed, $x$ has the weaker decay everywhere in the complete
gap. Intermediate angles were not computed.

% ---------------------------------------------------------------------
\section{What this means for the finite plate and the insertion loss}

\begin{enumerate}
  \item \emph{Transmission falls roughly like $e^{-\sigma L}$.} For a receiver at
        distance $L$, in-gap transmission drops by about $8.686\,\sigma L$ dB. The
        finite plate spans about ten cells ($1\,\mathrm{m}$); along $x$ at mid-gap that
        is about $35\,\mathrm{dB}$.
  \item \emph{The gap visible on a finite plate is narrower than the Bloch gap.} Near
        the edges $\sigma \to 0$ like a square root. For at least $20\,\mathrm{dB}$ over
        $1\,\mathrm{m}$ one needs $8.686\,\sigma L \ge 20$, i.e.\ $\sigma \ge
        2.3\,\mathrm{m}^{-1}$ or $\sigma a \ge 0.23$. Along $x$ this holds for roughly
        $880$--$1180\,\mathrm{Hz}$, compared with the Bloch gap of
        $1037$--$1200\,\mathrm{Hz}$ in this model.
  \item \emph{Evanescent decay is not damping.} It dissipates no energy: the lattice
        simply cannot carry energy away, so the energy stays near the source or is
        reflected back to it.
  \item \emph{Damping is a separate, weaker decay.} With a loss factor $\eta$ the
        bending wavenumber becomes $\kappa = \kappa_B(1 + i\eta)^{-1/4} \approx
        \kappa_B(1 - i\eta/4)$, so $\sigma = \eta\kappa_B/4$. At $1100\,\mathrm{Hz}$:
        $0.65\,\mathrm{dB/m}$ for $\eta = 0.01$ and $2.6\,\mathrm{dB/m}$ for
        $\eta = 0.04$. Inside the gap the $10$--$35\,\mathrm{dB/m}$ of the evanescent
        decay dominates, except right at the edges, which damping rounds off. Outside
        the gap damping is the only decay and lowers the resonance peaks of the finite
        plate. This is also why undamped Bloch bands suffice for the band positions.
  \item \emph{Few plate modes exist inside the gap.} A mode of the finite plate is a
        superposition of waves bouncing between its edges; inside the gap there is no
        propagating wave to bounce, so only motion localized at the source or at the
        edges remains.
  \item \emph{The driving point is not isolated.} The cell under the shaker still
        moves, through the near field and the decaying wave. The drive-point mobility
        drops much less than the transfer mobility far away. Which insertion-loss
        definition is used (mobility, radiated power) decides how strongly the gap
        shows; to be checked against the forced-response runs.
  \item \emph{Prerequisite:} all of this applies to a finite plate with resonators on
        the $0.1\,\mathrm{m}$ lattice (about 100 small resonators), not to the current
        layouts with 1--6 large resonators.
\end{enumerate}

% ---------------------------------------------------------------------
\section{How FE would compute $k(\omega)$}

In $K_r = T^H K T$ the wavevector enters only through the phase factors
$e^{\pm i k_x a}$ on the periodic images. With $k_y$ fixed and $z = e^{i k_x a}$, the
reduced problem takes the form
\begin{equation}
  \left(\frac{A_{-}}{z} + A_0 + A_{+} z\right)\psi = 0 ,
\end{equation}
with matrices $A$ that contain $\omega^2$: a \emph{quadratic eigenvalue problem} in $z$
at fixed $\omega$. Its roots come in pairs $z$ and $1/z$, i.e.\ $k$ and $-k$. For
complex $k$, $T^H$ has to be replaced by $T(-\mathbf{k})^{\mathsf{T}}$; conjugation is
the same thing only for real $k$. This is the $k(\omega)$ formulation listed as an
optional later step (damped gap depth); it is not needed for the band positions.

% ---------------------------------------------------------------------
\section{Check questions with answers}

\paragraph{Question 1.} \emph{Along $x$ the decaying root has $k_\mathrm{re} = \pi/a$
throughout the gap. What does the motion look like from one cell to the next, and why
must $\sigma$ go to zero at both gap edges?}

\emph{Answer.} Neighbouring cells move in antiphase ($e^{i\pi} = -1$), and the
amplitude shrinks by $e^{-\sigma a}$ per cell; at mid-gap, $\sigma a \approx 0.4$, so
each cell has about $0.67$ of the amplitude of the previous one. Each gap edge is the
flat extreme of a propagating band: the maximum of band 1 at $X$ below, the decoupled
sine wave at $X$ above. There $k(\omega)$ turns from real to complex, and $\sigma$
grows from zero like the square root of the distance from the edge.

\paragraph{Question 2.} \emph{The Bloch complete gap is $1031$--$1195\,\mathrm{Hz}$. On
the $1\,\mathrm{m}$ finite plate, in which part of that range would you expect a clear
insertion loss, and how does a loss factor of $0.04$ change the picture inside and
outside the gap?}

\emph{Answer.} The clearly visible range is narrower than the Bloch gap, because
$\sigma L$ is limited: near the edges $\sigma \to 0$, and $1\,\mathrm{m}$ is not long
enough to build up much attenuation. In the Kirchhoff model, at least $20\,\mathrm{dB}$
over $1\,\mathrm{m}$ along $x$ holds for roughly $880$--$1180\,\mathrm{Hz}$. The gap does
not ``dampen'': evanescent decay dissipates nothing, the lattice just cannot carry the
energy away. A loss factor of $0.04$ adds about $2.6\,\mathrm{dB/m}$ everywhere. Inside
the gap that is small against $10$--$35\,\mathrm{dB/m}$ and only rounds off the edges;
outside the gap it is the only decay and lowers the resonance peaks of the finite
plate, which also enters an insertion loss measured against the bare plate.

\paragraph{Question 3.} \emph{The diagonal decays about four times faster than $x$. A
receiver sits on the finite plate at some angle from the shaker. Which direction
controls how strongly the gap shows up at that receiver, and why?}

\emph{Answer.} Roughly the direction of the line from shaker to receiver relative to
the lattice axes. A receiver along $x$ or $y$ sees the weakest decay (at most about
$35\,\mathrm{dB/m}$), one along the diagonal about four times more. For a quantity
collected over the whole plate, such as radiated power or plate-averaged mobility, the
weakest direction dominates, since most of the energy that gets far travels along $x$
and $y$. Two caveats: the computed rates are for straight wavefronts along $x$ and the
diagonal, not the field of a point source, and on a finite plate edge reflections add
other paths. The faster diagonal decay is not ``four waves being damped'' (there is no
damping in this model, and the factor four is not a count). Two plausible contributions,
not separated: at $M$, four degenerate plate waves form one coupled combination with
$|\varphi|^2 = 4/(\rho h A)$ against $2/(\rho h A)$ at $X$, so the coupling to the
resonator is twice as strong; and $f_r$ lies inside the diagonal gap
($1037$--$1270\,\mathrm{Hz}$), where $\mu_\mathrm{eff}$ is large, while along $x$ it sits
at the upper edge, where the sine wave does not see the resonator at all.

% =====================================================================

\chapter{Inside a gap: complex wavenumbers and the finite plate}
\label{ch:complex_k}

The dispersion diagram shows where waves \emph{cannot} propagate. This chapter asks
what the plate does at a frequency inside a gap when a shaker forces it to move
there: the Bloch wavenumber becomes complex, the waves decay from cell to cell, and
the decay rate tells how strongly a finite plate of about ten cells can show the
gap. Numbers come from a Kirchhoff plane-wave model of the production cell
($a = 0.1\,\mathrm{m}$, $h = 5\,\mathrm{mm}$, steel, one resonator per cell,
$f_r = 1200\,\mathrm{Hz}$, $\mu = 0.25$), solved for complex $k$.

\begin{center}
\begin{tabular}{|l|p{0.68\textwidth}|}
\hline
Symbol & Meaning \\
\hline
$k$ & Bloch wavenumber along one direction; here allowed to be complex \\
$k_\mathrm{re}$, $\sigma$ & real and imaginary part of $k$; $\sigma$ is the decay rate in $1/\mathrm{m}$ \\
$\lambda = \omega^2$ & squared angular frequency \\
$\lambda_r$, $\lambda_X$ & resonator eigenvalue $(2\pi\cdot 1200\,\mathrm{Hz})^2$; bare plate eigenvalue at $X$ \\
$\mu$, $\beta = \lambda/\lambda_r$ & mass ratio $0.25$; squared frequency ratio \\
$\mu_\mathrm{eff} = \mu/(1-\beta)$ & effective mass ratio: the resonator acts on the plate like the mass $m/(1-\beta)$ (Chapter~\ref{ch:gap_position}) \\
$D$, $\rho h$ & bending stiffness and mass per area of the plate \\
$\kappa_B$ & bending wavenumber of the bare plate, $\kappa_B^4 = \omega^2\rho h/D$ \\
$p = \pi/a$ & wavenumber at the zone edge $X$ \\
$\delta$ & offset from the zone edge, $k = p + \delta$ \\
$\bar{H}$, $\Delta$, $c$ & mean and half difference of the diagonal, and off-diagonal entry, of the two-level matrix (Chapter~\ref{ch:eigenvectors}) \\
$s = \mathrm{d}\Delta/\mathrm{d}k$ & how fast $\Delta$ changes with $k$ \\
$\eta$ & loss factor of the plate \\
$L$ & distance from source to receiver \\
\hline
\end{tabular}
\end{center}

% ---------------------------------------------------------------------
\section{From $\omega(k)$ to $k(\omega)$}

The dispersion diagram answers: \emph{given a real $k$, which $\omega$ are
possible?} A forced response asks the reverse: the shaker fixes $\omega$, so
\emph{which $k$ are possible?} For real $k$, a frequency inside a gap has no answer.
That does not mean the plate cannot move at that frequency, since the shaker forces
it to. It means that only solutions with \emph{complex} $k$ exist there.

% ---------------------------------------------------------------------
\section{What a complex wavenumber means}

The Bloch theorem states $\psi(x + a) = e^{ika}\psi(x)$. With
$k = k_\mathrm{re} + i\sigma$,
\begin{equation}
  e^{ika} = e^{i k_\mathrm{re} a}\cdot e^{-\sigma a} .
\end{equation}
\begin{itemize}
  \item $e^{i k_\mathrm{re} a}$ is the phase step from one cell to the next, as before;
  \item $e^{-\sigma a}$ is the amplitude factor per cell: each cell multiplies the
        amplitude by $e^{-\sigma a}$.
\end{itemize}
In decibels, the decay per cell is
\begin{equation}
  20\log_{10}\!\left(e^{\sigma a}\right) = 20\,\sigma a\,\log_{10} e = 8.686\,\sigma a
  \ \mathrm{dB}, \qquad \text{or } 8.686\,\sigma \ \mathrm{dB/m} .
\end{equation}

Two special cases:
\begin{itemize}
  \item $k_\mathrm{re} = \pi/a$: $e^{i\pi} = -1$, so neighbouring cells move in
        \emph{antiphase} while the amplitude shrinks: a Bragg-type decaying standing
        pattern;
  \item $k_\mathrm{re} = 0$: neighbouring cells move \emph{in phase} while the amplitude
        shrinks.
\end{itemize}

\paragraph{Roots come in sets of four.} The plate equation contains $k$ only through
squares, so with $k$ also $-k$ is a solution. Without damping its coefficients are
real, so with $k$ also $k^*$ is a solution. The roots therefore come as
$\{k, -k, k^*, -k^*\}$. On each side of a source one keeps the root that decays away
from it.

% ---------------------------------------------------------------------
\section{Evanescent waves exist without resonators}

For the bare Kirchhoff plate, $\kappa^4 = \kappa_B^4$ has the four roots
$\kappa = \pm\kappa_B$ and $\pm i\kappa_B$. The imaginary pair is the \emph{bending
near field}, present at every frequency. At $1100\,\mathrm{Hz}$,
$\kappa_B = 30.1\,\mathrm{m}^{-1}$, so the near field decays by
$\kappa_B a = 3.0$ per cell ($26\,\mathrm{dB}$ per cell). It matters only within
about one cell of a source or an edge.

A gap is therefore not where evanescent waves begin. It is where \emph{the
propagating pair disappears}, and the slowest-decaying remaining wave is what is seen
at a distance.

% ---------------------------------------------------------------------
\section{Two-level model with complex $k$}
\label{sec:two_level_complex}

\subsection{The matrix near $X$}
Near $X$ along $x$, two plate waves matter: $e^{ikx}$ with
$\lambda_1^0 = (D/\rho h)\,k^4$, and its folded partner with
$\lambda_2^0 = (D/\rho h)\,(k - 2\pi/a)^4$. The resonator acts on the plate as a
frequency-dependent mass, so the rigid-mass matrix of Chapter~\ref{ch:two_level}
applies with $\mu$ replaced by $\mu_\mathrm{eff}$:
\begin{equation}
  H_{11} = \lambda_1^0 - \mu_\mathrm{eff}\lambda, \qquad
  H_{22} = \lambda_2^0 - \mu_\mathrm{eff}\lambda, \qquad
  c = -\mu_\mathrm{eff}\lambda .
\end{equation}
As in Chapter~\ref{ch:eigenvectors}, the eigenvalues are
$\lambda = \bar{H} \pm \sqrt{\Delta^2 + c^2}$.

\subsection{Step 1: $\Delta$ and $\bar{H}$ near the zone edge}
Write $k = p + \delta$ with $p = \pi/a$. Then $k - 2\pi/a = \delta - p$, and since
the fourth power is even,
\begin{equation}
  \lambda_1^0 = \frac{D}{\rho h}(p + \delta)^4, \qquad
  \lambda_2^0 = \frac{D}{\rho h}(p - \delta)^4 .
\end{equation}
Expanding both:
\begin{align}
  (p + \delta)^4 &= p^4 + 4p^3\delta + 6p^2\delta^2 + 4p\delta^3 + \delta^4, \\
  (p - \delta)^4 &= p^4 - 4p^3\delta + 6p^2\delta^2 - 4p\delta^3 + \delta^4 .
\end{align}
The difference keeps the odd powers, the sum the even ones:
\begin{align}
  \Delta &= \frac{\lambda_1^0 - \lambda_2^0}{2}
          = \frac{D}{\rho h}\left(4p^3\delta + 4p\delta^3\right)
          = \frac{4\lambda_X}{p}\,\delta\left(1 + \frac{\delta^2}{p^2}\right), \\
  \bar{H} &= \frac{\lambda_1^0 + \lambda_2^0}{2} - \mu_\mathrm{eff}\lambda
          = \lambda_X\left(1 + 6\frac{\delta^2}{p^2} + \frac{\delta^4}{p^4}\right) - \mu_\mathrm{eff}\lambda ,
\end{align}
using $\lambda_X = (D/\rho h)\,p^4$. To first order in $\delta$:
\begin{equation}
  \Delta \approx s\,\delta, \qquad s = \frac{4\lambda_X}{p} = \frac{4\lambda_X a}{\pi},
  \qquad \bar{H} \approx \lambda_X - \mu_\mathrm{eff}\lambda .
\end{equation}
The neglected term $6\delta^2/p^2$ in $\bar{H}$ is about $10\,\%$ in the middle of
this gap.

\subsection{Step 2: insert an imaginary offset}
Set $\delta = i\sigma$. Then $\Delta = is\sigma$ and $\Delta^2 = -s^2\sigma^2$, so
\begin{equation}
  \lambda = \bar{H} \pm \sqrt{c^2 - s^2\sigma^2} .
\end{equation}

\subsection{Step 3: solve for the decay rate}
Squaring, $(\lambda - \bar{H})^2 = c^2 - s^2\sigma^2$, hence
\begin{equation}
  \boxed{\;\sigma = \frac{\sqrt{c^2 - (\lambda - \bar{H})^2}}{s}\;}
  \label{eq:sigma_two_level}
\end{equation}
\begin{itemize}
  \item \emph{Inside the gap}, $|\lambda - \bar{H}| < |c|$: $\sigma$ is real, and
        $k = \pi/a + i\sigma$ is a decaying wave.
  \item \emph{Outside the gap}, the root would be imaginary, which makes $\delta$ real:
        an ordinary propagating wave.
  \item \emph{At an edge}, $\lambda - \bar{H} = \pm|c|$, so $\sigma = 0$. Let
        $\varepsilon$ be the distance of $\lambda$ from that edge. Then
        $c^2 - (\lambda - \bar{H})^2 = (|c| - |\lambda - \bar{H}|)(|c| + |\lambda - \bar{H}|)
        \approx \varepsilon\cdot 2|c|$, so
        \begin{equation}
          \sigma \approx \frac{\sqrt{2|c|\,\varepsilon}}{s} :
        \end{equation}
        $\sigma$ grows like the square root of the distance from the edge.
  \item \emph{With constant coupling}, $\sigma$ is largest at the gap centre:
        $\sigma_\mathrm{max} = |c|/s$, the gap half-width divided by how fast the two
        crossing bands separate in $k$.
\end{itemize}

\paragraph{The same from the band edge.} At a band edge the band is flat:
$\omega \approx \omega_e - \alpha(k - k_e)^2$ with $\alpha > 0$. Solving for $k$,
\begin{equation}
  k - k_e = \pm\sqrt{\frac{\omega_e - \omega}{\alpha}},
\end{equation}
which is real on the band side ($\omega < \omega_e$) and imaginary on the gap side.

\subsection{The resonator version}
Here $c = -\mu_\mathrm{eff}\lambda$ depends on the frequency and grows towards $f_r$.
The decoupled sine wave pins the upper edge at $\lambda_X$:
$\bar{H} = \lambda_X - |c|$. With $x = \lambda_X - \lambda$,
\begin{equation}
  \lambda - \bar{H} = |c| - x
  \quad\Rightarrow\quad
  c^2 - (\lambda - \bar{H})^2 = c^2 - (|c| - x)^2 = x\,(2|c| - x),
\end{equation}
so
\begin{equation}
  \sigma a = \frac{a\sqrt{x(2|c| - x)}}{s} = \frac{\pi}{4}\,\frac{\sqrt{x(2|c| - x)}}{\lambda_X} .
\end{equation}
Worked at $1050\,\mathrm{Hz}$, in units of $\lambda_X$ ($f_X = 1199.8\,\mathrm{Hz}$
in this model):
\begin{align}
  \lambda/\lambda_X &= (1050/1199.8)^2 = 0.7659, \qquad x = 0.2341, \\
  \beta &= (1050/1200)^2 = 0.7656, \qquad \mu_\mathrm{eff} = 0.25/0.2344 = 1.0666, \\
  |c| &= 1.0666\cdot 0.7659 = 0.8169, \\
  x(2|c| - x) &= 0.2341\cdot 1.3997 = 0.3277, \qquad \sigma a = \tfrac{\pi}{4}\sqrt{0.3277} = 0.450 .
\end{align}

\begin{center}
\begin{tabular}{|l|l|l|}
\hline
$f$ in Hz & $\sigma a$, two-wave formula & $\sigma a$, full plane-wave model \\
\hline
900 & 0.236 & 0.276 \\
1050 & 0.450 & 0.407 \\
1150 & 0.528 & 0.346 \\
\hline
\end{tabular}
\end{center}

The two-wave formula gets the size right and degrades towards $f_r$. There
$\mu_\mathrm{eff}$ becomes large, the resonator nearly pins its attachment point, and
many plane waves are needed to describe that, as in Chapter~\ref{ch:pwe_general}.

% ---------------------------------------------------------------------
\section{Results of the full plane-wave model}

\paragraph{Method.} The plane-wave condition of Chapter~\ref{ch:pwe_general},
\begin{equation}
  1 = \mu_\mathrm{eff}\sum_{\mathbf{G}} \frac{\lambda}{\lambda_\mathbf{G}(\mathbf{k}) - \lambda},
  \qquad
  \lambda_\mathbf{G}(\mathbf{k}) = \frac{D}{\rho h}\left[(k_x + G_x)^2 + (k_y + G_y)^2\right]^2,
\end{equation}
with Kirchhoff plate waves and $31\times 31 = 961$ harmonics, is solved at fixed
frequency for complex $k$ along a direction. The square $(k_x + G_x)^2$ is written
without absolute value, so it continues to complex $k$. Roots are located from minima
of the residual on a grid in the complex plane and polished with Newton's method.

\paragraph{Checks.} At real $k$ the model gives the band frequencies $839.2$ and
$1660.6\,\mathrm{Hz}$ at $X$ and $1036.7\,\mathrm{Hz}$ at $M$ (FE: $836$, $1653$,
$1031$). At a pass-band frequency it returns the known real $k$. The complete gap of
this model is $1036.7$--$1199.8\,\mathrm{Hz}$ (FE: $1031$--$1195\,\mathrm{Hz}$).

\paragraph{Along $x$ ($\Gamma$--$X$), slowest-decaying root:}

\begin{center}
\begin{tabular}{|l|l|l|l|}
\hline
$f$ in Hz & $k_\mathrm{re}/(\pi/a)$ & $\sigma a$ & dB/m \\
\hline
800 & 0.921 & 0 & 0 (propagating) \\
841 & 1 & 0.051 & 4.5 \\
850 & 1 & 0.124 & 10.7 \\
900 & 1 & 0.276 & 24.0 \\
1050 & 1 & 0.407 & 35.4 (maximum) \\
1150 & 1 & 0.346 & 30.1 \\
1190 & 1 & 0.191 & 16.6 \\
1198 & 1 & 0.087 & 7.5 \\
\hline
\end{tabular}
\end{center}

\begin{itemize}
  \item The directional gap along $x$ is $839$--$1200\,\mathrm{Hz}$, and the decaying
        root stays at $k_\mathrm{re} = \pi/a$ throughout: antiphase from cell to cell.
  \item Square-root check at the lower edge: $\sigma a/\sqrt{f - 839.2\,\mathrm{Hz}} =
        0.038$ at $841$, $845$ and $850\,\mathrm{Hz}$, constant as predicted.
  \item The near-field root ($k_\mathrm{re} = 0$, $\sigma a \approx 2.9$--$3.7$) is
        always present and plays no role at a distance.
\end{itemize}

\paragraph{Along the diagonal ($\Gamma$--$M$):}

\begin{center}
\begin{tabular}{|l|l|l|l|}
\hline
$f$ in Hz & $k_\mathrm{re}/(\sqrt{2}\pi/a)$ & $\sigma a$ & dB/m \\
\hline
1000 & 0.816 & 0 & 0 \\
1040 & 1 & 0.289 & 25 \\
1050 & 1 & 0.615 & 53 \\
1075 & 0.934 & 1.463 & 127 \\
1150 & 0.697 & 1.532 & 133 \\
1210 & 0.561 & 1.387 & 120 \\
1265 & 0.386 & 0.701 & 61 \\
1300 & 0.549 & 0 & 0 \\
\hline
\end{tabular}
\end{center}

\begin{itemize}
  \item Between $1050$ and $1075\,\mathrm{Hz}$, two roots at the zone edge meet and
        split into a complex pair; after that $k_\mathrm{re}$ leaves the zone edge, and
        the wave oscillates while it decays.
  \item The diagonal gap reaches up to about $1270\,\mathrm{Hz}$ and decays roughly
        four times faster than along $x$.
\end{itemize}

Of the two directions computed, $x$ has the weaker decay everywhere in the complete
gap. Intermediate angles were not computed.

% ---------------------------------------------------------------------
\section{What this means for the finite plate and the insertion loss}

\begin{enumerate}
  \item \emph{Transmission falls roughly like $e^{-\sigma L}$.} For a receiver at
        distance $L$, in-gap transmission drops by about $8.686\,\sigma L$ dB. The
        finite plate spans about ten cells ($1\,\mathrm{m}$); along $x$ at mid-gap that
        is about $35\,\mathrm{dB}$.
  \item \emph{The gap visible on a finite plate is narrower than the Bloch gap.} Near
        the edges $\sigma \to 0$ like a square root. For at least $20\,\mathrm{dB}$ over
        $1\,\mathrm{m}$ one needs $8.686\,\sigma L \ge 20$, i.e.\ $\sigma \ge
        2.3\,\mathrm{m}^{-1}$ or $\sigma a \ge 0.23$. Along $x$ this holds for roughly
        $880$--$1180\,\mathrm{Hz}$, compared with the Bloch gap of
        $1037$--$1200\,\mathrm{Hz}$ in this model.
  \item \emph{Evanescent decay is not damping.} It dissipates no energy: the lattice
        simply cannot carry energy away, so the energy stays near the source or is
        reflected back to it.
  \item \emph{Damping is a separate, weaker decay.} With a loss factor $\eta$ the
        bending wavenumber becomes $\kappa = \kappa_B(1 + i\eta)^{-1/4} \approx
        \kappa_B(1 - i\eta/4)$, so $\sigma = \eta\kappa_B/4$. At $1100\,\mathrm{Hz}$:
        $0.65\,\mathrm{dB/m}$ for $\eta = 0.01$ and $2.6\,\mathrm{dB/m}$ for
        $\eta = 0.04$. Inside the gap the $10$--$35\,\mathrm{dB/m}$ of the evanescent
        decay dominates, except right at the edges, which damping rounds off. Outside
        the gap damping is the only decay and lowers the resonance peaks of the finite
        plate. This is also why undamped Bloch bands suffice for the band positions.
  \item \emph{Few plate modes exist inside the gap.} A mode of the finite plate is a
        superposition of waves bouncing between its edges; inside the gap there is no
        propagating wave to bounce, so only motion localized at the source or at the
        edges remains.
  \item \emph{The driving point is not isolated.} The cell under the shaker still
        moves, through the near field and the decaying wave. The drive-point mobility
        drops much less than the transfer mobility far away. Which insertion-loss
        definition is used (mobility, radiated power) decides how strongly the gap
        shows; to be checked against the forced-response runs.
  \item \emph{Prerequisite:} all of this applies to a finite plate with resonators on
        the $0.1\,\mathrm{m}$ lattice (about 100 small resonators), not to the current
        layouts with 1--6 large resonators.
\end{enumerate}

% ---------------------------------------------------------------------
\section{How FE would compute $k(\omega)$}

In $K_r = T^H K T$ the wavevector enters only through the phase factors
$e^{\pm i k_x a}$ on the periodic images. With $k_y$ fixed and $z = e^{i k_x a}$, the
reduced problem takes the form
\begin{equation}
  \left(\frac{A_{-}}{z} + A_0 + A_{+} z\right)\psi = 0 ,
\end{equation}
with matrices $A$ that contain $\omega^2$: a \emph{quadratic eigenvalue problem} in $z$
at fixed $\omega$. Its roots come in pairs $z$ and $1/z$, i.e.\ $k$ and $-k$. For
complex $k$, $T^H$ has to be replaced by $T(-\mathbf{k})^{\mathsf{T}}$; conjugation is
the same thing only for real $k$. This is the $k(\omega)$ formulation listed as an
optional later step (damped gap depth); it is not needed for the band positions.

% ---------------------------------------------------------------------
\section{Check questions with answers}

\paragraph{Question 1.} \emph{Along $x$ the decaying root has $k_\mathrm{re} = \pi/a$
throughout the gap. What does the motion look like from one cell to the next, and why
must $\sigma$ go to zero at both gap edges?}

\emph{Answer.} Neighbouring cells move in antiphase ($e^{i\pi} = -1$), and the
amplitude shrinks by $e^{-\sigma a}$ per cell; at mid-gap, $\sigma a \approx 0.4$, so
each cell has about $0.67$ of the amplitude of the previous one. Each gap edge is the
flat extreme of a propagating band: the maximum of band 1 at $X$ below, the decoupled
sine wave at $X$ above. There $k(\omega)$ turns from real to complex, and $\sigma$
grows from zero like the square root of the distance from the edge.

\paragraph{Question 2.} \emph{The Bloch complete gap is $1031$--$1195\,\mathrm{Hz}$. On
the $1\,\mathrm{m}$ finite plate, in which part of that range would you expect a clear
insertion loss, and how does a loss factor of $0.04$ change the picture inside and
outside the gap?}

\emph{Answer.} The clearly visible range is narrower than the Bloch gap, because
$\sigma L$ is limited: near the edges $\sigma \to 0$, and $1\,\mathrm{m}$ is not long
enough to build up much attenuation. In the Kirchhoff model, at least $20\,\mathrm{dB}$
over $1\,\mathrm{m}$ along $x$ holds for roughly $880$--$1180\,\mathrm{Hz}$. The gap does
not ``dampen'': evanescent decay dissipates nothing, the lattice just cannot carry the
energy away. A loss factor of $0.04$ adds about $2.6\,\mathrm{dB/m}$ everywhere. Inside
the gap that is small against $10$--$35\,\mathrm{dB/m}$ and only rounds off the edges;
outside the gap it is the only decay and lowers the resonance peaks of the finite
plate, which also enters an insertion loss measured against the bare plate.

\paragraph{Question 3.} \emph{The diagonal decays about four times faster than $x$. A
receiver sits on the finite plate at some angle from the shaker. Which direction
controls how strongly the gap shows up at that receiver, and why?}

\emph{Answer.} Roughly the direction of the line from shaker to receiver relative to
the lattice axes. A receiver along $x$ or $y$ sees the weakest decay (at most about
$35\,\mathrm{dB/m}$), one along the diagonal about four times more. For a quantity
collected over the whole plate, such as radiated power or plate-averaged mobility, the
weakest direction dominates, since most of the energy that gets far travels along $x$
and $y$. Two caveats: the computed rates are for straight wavefronts along $x$ and the
diagonal, not the field of a point source, and on a finite plate edge reflections add
other paths. The faster diagonal decay is not ``four waves being damped'' (there is no
damping in this model, and the factor four is not a count). Two plausible contributions,
not separated: at $M$, four degenerate plate waves form one coupled combination with
$|\varphi|^2 = 4/(\rho h A)$ against $2/(\rho h A)$ at $X$, so the coupling to the
resonator is twice as strong; and $f_r$ lies inside the diagonal gap
($1037$--$1270\,\mathrm{Hz}$), where $\mu_\mathrm{eff}$ is large, while along $x$ it sits
at the upper edge, where the sine wave does not see the resonator at all.

% =====================================================================

\chapter{Inside a gap: complex wavenumbers and the finite plate}
\label{ch:complex_k}

The dispersion diagram shows where waves \emph{cannot} propagate. This chapter asks
what the plate does at a frequency inside a gap when a shaker forces it to move
there: the Bloch wavenumber becomes complex, the waves decay from cell to cell, and
the decay rate tells how strongly a finite plate of about ten cells can show the
gap. Numbers come from a Kirchhoff plane-wave model of the production cell
($a = 0.1\,\mathrm{m}$, $h = 5\,\mathrm{mm}$, steel, one resonator per cell,
$f_r = 1200\,\mathrm{Hz}$, $\mu = 0.25$), solved for complex $k$.

\begin{center}
\begin{tabular}{|l|p{0.68\textwidth}|}
\hline
Symbol & Meaning \\
\hline
$k$ & Bloch wavenumber along one direction; here allowed to be complex \\
$k_\mathrm{re}$, $\sigma$ & real and imaginary part of $k$; $\sigma$ is the decay rate in $1/\mathrm{m}$ \\
$\lambda = \omega^2$ & squared angular frequency \\
$\lambda_r$, $\lambda_X$ & resonator eigenvalue $(2\pi\cdot 1200\,\mathrm{Hz})^2$; bare plate eigenvalue at $X$ \\
$\mu$, $\beta = \lambda/\lambda_r$ & mass ratio $0.25$; squared frequency ratio \\
$\mu_\mathrm{eff} = \mu/(1-\beta)$ & effective mass ratio: the resonator acts on the plate like the mass $m/(1-\beta)$ (Chapter~\ref{ch:gap_position}) \\
$D$, $\rho h$ & bending stiffness and mass per area of the plate \\
$\kappa_B$ & bending wavenumber of the bare plate, $\kappa_B^4 = \omega^2\rho h/D$ \\
$p = \pi/a$ & wavenumber at the zone edge $X$ \\
$\delta$ & offset from the zone edge, $k = p + \delta$ \\
$\bar{H}$, $\Delta$, $c$ & mean and half difference of the diagonal, and off-diagonal entry, of the two-level matrix (Chapter~\ref{ch:eigenvectors}) \\
$s = \mathrm{d}\Delta/\mathrm{d}k$ & how fast $\Delta$ changes with $k$ \\
$\eta$ & loss factor of the plate \\
$L$ & distance from source to receiver \\
\hline
\end{tabular}
\end{center}

% ---------------------------------------------------------------------
\section{From $\omega(k)$ to $k(\omega)$}

The dispersion diagram answers: \emph{given a real $k$, which $\omega$ are
possible?} A forced response asks the reverse: the shaker fixes $\omega$, so
\emph{which $k$ are possible?} For real $k$, a frequency inside a gap has no answer.
That does not mean the plate cannot move at that frequency, since the shaker forces
it to. It means that only solutions with \emph{complex} $k$ exist there.

% ---------------------------------------------------------------------
\section{What a complex wavenumber means}

The Bloch theorem states $\psi(x + a) = e^{ika}\psi(x)$. With
$k = k_\mathrm{re} + i\sigma$,
\begin{equation}
  e^{ika} = e^{i k_\mathrm{re} a}\cdot e^{-\sigma a} .
\end{equation}
\begin{itemize}
  \item $e^{i k_\mathrm{re} a}$ is the phase step from one cell to the next, as before;
  \item $e^{-\sigma a}$ is the amplitude factor per cell: each cell multiplies the
        amplitude by $e^{-\sigma a}$.
\end{itemize}
In decibels, the decay per cell is
\begin{equation}
  20\log_{10}\!\left(e^{\sigma a}\right) = 20\,\sigma a\,\log_{10} e = 8.686\,\sigma a
  \ \mathrm{dB}, \qquad \text{or } 8.686\,\sigma \ \mathrm{dB/m} .
\end{equation}

Two special cases:
\begin{itemize}
  \item $k_\mathrm{re} = \pi/a$: $e^{i\pi} = -1$, so neighbouring cells move in
        \emph{antiphase} while the amplitude shrinks: a Bragg-type decaying standing
        pattern;
  \item $k_\mathrm{re} = 0$: neighbouring cells move \emph{in phase} while the amplitude
        shrinks.
\end{itemize}

\paragraph{Roots come in sets of four.} The plate equation contains $k$ only through
squares, so with $k$ also $-k$ is a solution. Without damping its coefficients are
real, so with $k$ also $k^*$ is a solution. The roots therefore come as
$\{k, -k, k^*, -k^*\}$. On each side of a source one keeps the root that decays away
from it.

% ---------------------------------------------------------------------
\section{Evanescent waves exist without resonators}

For the bare Kirchhoff plate, $\kappa^4 = \kappa_B^4$ has the four roots
$\kappa = \pm\kappa_B$ and $\pm i\kappa_B$. The imaginary pair is the \emph{bending
near field}, present at every frequency. At $1100\,\mathrm{Hz}$,
$\kappa_B = 30.1\,\mathrm{m}^{-1}$, so the near field decays by
$\kappa_B a = 3.0$ per cell ($26\,\mathrm{dB}$ per cell). It matters only within
about one cell of a source or an edge.

A gap is therefore not where evanescent waves begin. It is where \emph{the
propagating pair disappears}, and the slowest-decaying remaining wave is what is seen
at a distance.

% ---------------------------------------------------------------------
\section{Two-level model with complex $k$}
\label{sec:two_level_complex}

\subsection{The matrix near $X$}
Near $X$ along $x$, two plate waves matter: $e^{ikx}$ with
$\lambda_1^0 = (D/\rho h)\,k^4$, and its folded partner with
$\lambda_2^0 = (D/\rho h)\,(k - 2\pi/a)^4$. The resonator acts on the plate as a
frequency-dependent mass, so the rigid-mass matrix of Chapter~\ref{ch:two_level}
applies with $\mu$ replaced by $\mu_\mathrm{eff}$:
\begin{equation}
  H_{11} = \lambda_1^0 - \mu_\mathrm{eff}\lambda, \qquad
  H_{22} = \lambda_2^0 - \mu_\mathrm{eff}\lambda, \qquad
  c = -\mu_\mathrm{eff}\lambda .
\end{equation}
As in Chapter~\ref{ch:eigenvectors}, the eigenvalues are
$\lambda = \bar{H} \pm \sqrt{\Delta^2 + c^2}$.

\subsection{Step 1: $\Delta$ and $\bar{H}$ near the zone edge}
Write $k = p + \delta$ with $p = \pi/a$. Then $k - 2\pi/a = \delta - p$, and since
the fourth power is even,
\begin{equation}
  \lambda_1^0 = \frac{D}{\rho h}(p + \delta)^4, \qquad
  \lambda_2^0 = \frac{D}{\rho h}(p - \delta)^4 .
\end{equation}
Expanding both:
\begin{align}
  (p + \delta)^4 &= p^4 + 4p^3\delta + 6p^2\delta^2 + 4p\delta^3 + \delta^4, \\
  (p - \delta)^4 &= p^4 - 4p^3\delta + 6p^2\delta^2 - 4p\delta^3 + \delta^4 .
\end{align}
The difference keeps the odd powers, the sum the even ones:
\begin{align}
  \Delta &= \frac{\lambda_1^0 - \lambda_2^0}{2}
          = \frac{D}{\rho h}\left(4p^3\delta + 4p\delta^3\right)
          = \frac{4\lambda_X}{p}\,\delta\left(1 + \frac{\delta^2}{p^2}\right), \\
  \bar{H} &= \frac{\lambda_1^0 + \lambda_2^0}{2} - \mu_\mathrm{eff}\lambda
          = \lambda_X\left(1 + 6\frac{\delta^2}{p^2} + \frac{\delta^4}{p^4}\right) - \mu_\mathrm{eff}\lambda ,
\end{align}
using $\lambda_X = (D/\rho h)\,p^4$. To first order in $\delta$:
\begin{equation}
  \Delta \approx s\,\delta, \qquad s = \frac{4\lambda_X}{p} = \frac{4\lambda_X a}{\pi},
  \qquad \bar{H} \approx \lambda_X - \mu_\mathrm{eff}\lambda .
\end{equation}
The neglected term $6\delta^2/p^2$ in $\bar{H}$ is about $10\,\%$ in the middle of
this gap.

\subsection{Step 2: insert an imaginary offset}
Set $\delta = i\sigma$. Then $\Delta = is\sigma$ and $\Delta^2 = -s^2\sigma^2$, so
\begin{equation}
  \lambda = \bar{H} \pm \sqrt{c^2 - s^2\sigma^2} .
\end{equation}

\subsection{Step 3: solve for the decay rate}
Squaring, $(\lambda - \bar{H})^2 = c^2 - s^2\sigma^2$, hence
\begin{equation}
  \boxed{\;\sigma = \frac{\sqrt{c^2 - (\lambda - \bar{H})^2}}{s}\;}
  \label{eq:sigma_two_level}
\end{equation}
\begin{itemize}
  \item \emph{Inside the gap}, $|\lambda - \bar{H}| < |c|$: $\sigma$ is real, and
        $k = \pi/a + i\sigma$ is a decaying wave.
  \item \emph{Outside the gap}, the root would be imaginary, which makes $\delta$ real:
        an ordinary propagating wave.
  \item \emph{At an edge}, $\lambda - \bar{H} = \pm|c|$, so $\sigma = 0$. Let
        $\varepsilon$ be the distance of $\lambda$ from that edge. Then
        $c^2 - (\lambda - \bar{H})^2 = (|c| - |\lambda - \bar{H}|)(|c| + |\lambda - \bar{H}|)
        \approx \varepsilon\cdot 2|c|$, so
        \begin{equation}
          \sigma \approx \frac{\sqrt{2|c|\,\varepsilon}}{s} :
        \end{equation}
        $\sigma$ grows like the square root of the distance from the edge.
  \item \emph{With constant coupling}, $\sigma$ is largest at the gap centre:
        $\sigma_\mathrm{max} = |c|/s$, the gap half-width divided by how fast the two
        crossing bands separate in $k$.
\end{itemize}

\paragraph{The same from the band edge.} At a band edge the band is flat:
$\omega \approx \omega_e - \alpha(k - k_e)^2$ with $\alpha > 0$. Solving for $k$,
\begin{equation}
  k - k_e = \pm\sqrt{\frac{\omega_e - \omega}{\alpha}},
\end{equation}
which is real on the band side ($\omega < \omega_e$) and imaginary on the gap side.

\subsection{The resonator version}
Here $c = -\mu_\mathrm{eff}\lambda$ depends on the frequency and grows towards $f_r$.
The decoupled sine wave pins the upper edge at $\lambda_X$:
$\bar{H} = \lambda_X - |c|$. With $x = \lambda_X - \lambda$,
\begin{equation}
  \lambda - \bar{H} = |c| - x
  \quad\Rightarrow\quad
  c^2 - (\lambda - \bar{H})^2 = c^2 - (|c| - x)^2 = x\,(2|c| - x),
\end{equation}
so
\begin{equation}
  \sigma a = \frac{a\sqrt{x(2|c| - x)}}{s} = \frac{\pi}{4}\,\frac{\sqrt{x(2|c| - x)}}{\lambda_X} .
\end{equation}
Worked at $1050\,\mathrm{Hz}$, in units of $\lambda_X$ ($f_X = 1199.8\,\mathrm{Hz}$
in this model):
\begin{align}
  \lambda/\lambda_X &= (1050/1199.8)^2 = 0.7659, \qquad x = 0.2341, \\
  \beta &= (1050/1200)^2 = 0.7656, \qquad \mu_\mathrm{eff} = 0.25/0.2344 = 1.0666, \\
  |c| &= 1.0666\cdot 0.7659 = 0.8169, \\
  x(2|c| - x) &= 0.2341\cdot 1.3997 = 0.3277, \qquad \sigma a = \tfrac{\pi}{4}\sqrt{0.3277} = 0.450 .
\end{align}

\begin{center}
\begin{tabular}{|l|l|l|}
\hline
$f$ in Hz & $\sigma a$, two-wave formula & $\sigma a$, full plane-wave model \\
\hline
900 & 0.236 & 0.276 \\
1050 & 0.450 & 0.407 \\
1150 & 0.528 & 0.346 \\
\hline
\end{tabular}
\end{center}

The two-wave formula gets the size right and degrades towards $f_r$. There
$\mu_\mathrm{eff}$ becomes large, the resonator nearly pins its attachment point, and
many plane waves are needed to describe that, as in Chapter~\ref{ch:pwe_general}.

% ---------------------------------------------------------------------
\section{Results of the full plane-wave model}

\paragraph{Method.} The plane-wave condition of Chapter~\ref{ch:pwe_general},
\begin{equation}
  1 = \mu_\mathrm{eff}\sum_{\mathbf{G}} \frac{\lambda}{\lambda_\mathbf{G}(\mathbf{k}) - \lambda},
  \qquad
  \lambda_\mathbf{G}(\mathbf{k}) = \frac{D}{\rho h}\left[(k_x + G_x)^2 + (k_y + G_y)^2\right]^2,
\end{equation}
with Kirchhoff plate waves and $31\times 31 = 961$ harmonics, is solved at fixed
frequency for complex $k$ along a direction. The square $(k_x + G_x)^2$ is written
without absolute value, so it continues to complex $k$. Roots are located from minima
of the residual on a grid in the complex plane and polished with Newton's method.

\paragraph{Checks.} At real $k$ the model gives the band frequencies $839.2$ and
$1660.6\,\mathrm{Hz}$ at $X$ and $1036.7\,\mathrm{Hz}$ at $M$ (FE: $836$, $1653$,
$1031$). At a pass-band frequency it returns the known real $k$. The complete gap of
this model is $1036.7$--$1199.8\,\mathrm{Hz}$ (FE: $1031$--$1195\,\mathrm{Hz}$).

\paragraph{Along $x$ ($\Gamma$--$X$), slowest-decaying root:}

\begin{center}
\begin{tabular}{|l|l|l|l|}
\hline
$f$ in Hz & $k_\mathrm{re}/(\pi/a)$ & $\sigma a$ & dB/m \\
\hline
800 & 0.921 & 0 & 0 (propagating) \\
841 & 1 & 0.051 & 4.5 \\
850 & 1 & 0.124 & 10.7 \\
900 & 1 & 0.276 & 24.0 \\
1050 & 1 & 0.407 & 35.4 (maximum) \\
1150 & 1 & 0.346 & 30.1 \\
1190 & 1 & 0.191 & 16.6 \\
1198 & 1 & 0.087 & 7.5 \\
\hline
\end{tabular}
\end{center}

\begin{itemize}
  \item The directional gap along $x$ is $839$--$1200\,\mathrm{Hz}$, and the decaying
        root stays at $k_\mathrm{re} = \pi/a$ throughout: antiphase from cell to cell.
  \item Square-root check at the lower edge: $\sigma a/\sqrt{f - 839.2\,\mathrm{Hz}} =
        0.038$ at $841$, $845$ and $850\,\mathrm{Hz}$, constant as predicted.
  \item The near-field root ($k_\mathrm{re} = 0$, $\sigma a \approx 2.9$--$3.7$) is
        always present and plays no role at a distance.
\end{itemize}

\paragraph{Along the diagonal ($\Gamma$--$M$):}

\begin{center}
\begin{tabular}{|l|l|l|l|}
\hline
$f$ in Hz & $k_\mathrm{re}/(\sqrt{2}\pi/a)$ & $\sigma a$ & dB/m \\
\hline
1000 & 0.816 & 0 & 0 \\
1040 & 1 & 0.289 & 25 \\
1050 & 1 & 0.615 & 53 \\
1075 & 0.934 & 1.463 & 127 \\
1150 & 0.697 & 1.532 & 133 \\
1210 & 0.561 & 1.387 & 120 \\
1265 & 0.386 & 0.701 & 61 \\
1300 & 0.549 & 0 & 0 \\
\hline
\end{tabular}
\end{center}

\begin{itemize}
  \item Between $1050$ and $1075\,\mathrm{Hz}$, two roots at the zone edge meet and
        split into a complex pair; after that $k_\mathrm{re}$ leaves the zone edge, and
        the wave oscillates while it decays.
  \item The diagonal gap reaches up to about $1270\,\mathrm{Hz}$ and decays roughly
        four times faster than along $x$.
\end{itemize}

Of the two directions computed, $x$ has the weaker decay everywhere in the complete
gap. Intermediate angles were not computed.

% ---------------------------------------------------------------------
\section{What this means for the finite plate and the insertion loss}

\begin{enumerate}
  \item \emph{Transmission falls roughly like $e^{-\sigma L}$.} For a receiver at
        distance $L$, in-gap transmission drops by about $8.686\,\sigma L$ dB. The
        finite plate spans about ten cells ($1\,\mathrm{m}$); along $x$ at mid-gap that
        is about $35\,\mathrm{dB}$.
  \item \emph{The gap visible on a finite plate is narrower than the Bloch gap.} Near
        the edges $\sigma \to 0$ like a square root. For at least $20\,\mathrm{dB}$ over
        $1\,\mathrm{m}$ one needs $8.686\,\sigma L \ge 20$, i.e.\ $\sigma \ge
        2.3\,\mathrm{m}^{-1}$ or $\sigma a \ge 0.23$. Along $x$ this holds for roughly
        $880$--$1180\,\mathrm{Hz}$, compared with the Bloch gap of
        $1037$--$1200\,\mathrm{Hz}$ in this model.
  \item \emph{Evanescent decay is not damping.} It dissipates no energy: the lattice
        simply cannot carry energy away, so the energy stays near the source or is
        reflected back to it.
  \item \emph{Damping is a separate, weaker decay.} With a loss factor $\eta$ the
        bending wavenumber becomes $\kappa = \kappa_B(1 + i\eta)^{-1/4} \approx
        \kappa_B(1 - i\eta/4)$, so $\sigma = \eta\kappa_B/4$. At $1100\,\mathrm{Hz}$:
        $0.65\,\mathrm{dB/m}$ for $\eta = 0.01$ and $2.6\,\mathrm{dB/m}$ for
        $\eta = 0.04$. Inside the gap the $10$--$35\,\mathrm{dB/m}$ of the evanescent
        decay dominates, except right at the edges, which damping rounds off. Outside
        the gap damping is the only decay and lowers the resonance peaks of the finite
        plate. This is also why undamped Bloch bands suffice for the band positions.
  \item \emph{Few plate modes exist inside the gap.} A mode of the finite plate is a
        superposition of waves bouncing between its edges; inside the gap there is no
        propagating wave to bounce, so only motion localized at the source or at the
        edges remains.
  \item \emph{The driving point is not isolated.} The cell under the shaker still
        moves, through the near field and the decaying wave. The drive-point mobility
        drops much less than the transfer mobility far away. Which insertion-loss
        definition is used (mobility, radiated power) decides how strongly the gap
        shows; to be checked against the forced-response runs.
  \item \emph{Prerequisite:} all of this applies to a finite plate with resonators on
        the $0.1\,\mathrm{m}$ lattice (about 100 small resonators), not to the current
        layouts with 1--6 large resonators.
\end{enumerate}

% ---------------------------------------------------------------------
\section{How FE would compute $k(\omega)$}

In $K_r = T^H K T$ the wavevector enters only through the phase factors
$e^{\pm i k_x a}$ on the periodic images. With $k_y$ fixed and $z = e^{i k_x a}$, the
reduced problem takes the form
\begin{equation}
  \left(\frac{A_{-}}{z} + A_0 + A_{+} z\right)\psi = 0 ,
\end{equation}
with matrices $A$ that contain $\omega^2$: a \emph{quadratic eigenvalue problem} in $z$
at fixed $\omega$. Its roots come in pairs $z$ and $1/z$, i.e.\ $k$ and $-k$. For
complex $k$, $T^H$ has to be replaced by $T(-\mathbf{k})^{\mathsf{T}}$; conjugation is
the same thing only for real $k$. This is the $k(\omega)$ formulation listed as an
optional later step (damped gap depth); it is not needed for the band positions.

% ---------------------------------------------------------------------
\section{Check questions with answers}

\paragraph{Question 1.} \emph{Along $x$ the decaying root has $k_\mathrm{re} = \pi/a$
throughout the gap. What does the motion look like from one cell to the next, and why
must $\sigma$ go to zero at both gap edges?}

\emph{Answer.} Neighbouring cells move in antiphase ($e^{i\pi} = -1$), and the
amplitude shrinks by $e^{-\sigma a}$ per cell; at mid-gap, $\sigma a \approx 0.4$, so
each cell has about $0.67$ of the amplitude of the previous one. Each gap edge is the
flat extreme of a propagating band: the maximum of band 1 at $X$ below, the decoupled
sine wave at $X$ above. There $k(\omega)$ turns from real to complex, and $\sigma$
grows from zero like the square root of the distance from the edge.

\paragraph{Question 2.} \emph{The Bloch complete gap is $1031$--$1195\,\mathrm{Hz}$. On
the $1\,\mathrm{m}$ finite plate, in which part of that range would you expect a clear
insertion loss, and how does a loss factor of $0.04$ change the picture inside and
outside the gap?}

\emph{Answer.} The clearly visible range is narrower than the Bloch gap, because
$\sigma L$ is limited: near the edges $\sigma \to 0$, and $1\,\mathrm{m}$ is not long
enough to build up much attenuation. In the Kirchhoff model, at least $20\,\mathrm{dB}$
over $1\,\mathrm{m}$ along $x$ holds for roughly $880$--$1180\,\mathrm{Hz}$. The gap does
not ``dampen'': evanescent decay dissipates nothing, the lattice just cannot carry the
energy away. A loss factor of $0.04$ adds about $2.6\,\mathrm{dB/m}$ everywhere. Inside
the gap that is small against $10$--$35\,\mathrm{dB/m}$ and only rounds off the edges;
outside the gap it is the only decay and lowers the resonance peaks of the finite
plate, which also enters an insertion loss measured against the bare plate.

\paragraph{Question 3.} \emph{The diagonal decays about four times faster than $x$. A
receiver sits on the finite plate at some angle from the shaker. Which direction
controls how strongly the gap shows up at that receiver, and why?}

\emph{Answer.} Roughly the direction of the line from shaker to receiver relative to
the lattice axes. A receiver along $x$ or $y$ sees the weakest decay (at most about
$35\,\mathrm{dB/m}$), one along the diagonal about four times more. For a quantity
collected over the whole plate, such as radiated power or plate-averaged mobility, the
weakest direction dominates, since most of the energy that gets far travels along $x$
and $y$. Two caveats: the computed rates are for straight wavefronts along $x$ and the
diagonal, not the field of a point source, and on a finite plate edge reflections add
other paths. The faster diagonal decay is not ``four waves being damped'' (there is no
damping in this model, and the factor four is not a count). Two plausible contributions,
not separated: at $M$, four degenerate plate waves form one coupled combination with
$|\varphi|^2 = 4/(\rho h A)$ against $2/(\rho h A)$ at $X$, so the coupling to the
resonator is twice as strong; and $f_r$ lies inside the diagonal gap
($1037$--$1270\,\mathrm{Hz}$), where $\mu_\mathrm{eff}$ is large, while along $x$ it sits
at the upper edge, where the sine wave does not see the resonator at all.
